\section{Uniform compilation and completeness}\label{uniform:section}

The algebraic witnesses are existential and are never supplied to the
graph compiler. We now enumerate only the product schemas, local blocks,
constant-size valuations, and graph vertices. The two realizations have
different output costs, so we retain their quantitative statements
separately. All operations are counted on deterministic multitape
Turing machines with sequential list access.

For the instance compiled from a fixed machine $M$ and polynomial $p$
as in \Cref{prop:computation}, on a word of length $m$, put
$r=2(p(m)+2)$ and $b_*=r+2$. Since $p(m)\ge m+2$, we have $m\le r$.
The common size calculation starts with $O_M(b_*^2)$ real types and
comparison schemas. Subdivision
multiplies this count by $O(r)$, and choosing a row for each local
block multiplies it by another $O(r)$. Thus there are $O_M(b_*^4)$
blocks, each with an absolute constant number of valuations.
The realizations then have the following growth:
\begin{center}
\begin{tabular}{@{}lll@{}}
\toprule
 & General realization & Subcubic realization\\
\midrule
Intermediate vertices & $O_M(b_*^4)$ base vertices
 & $O_M(b_*^8)$ auxiliary vertices\\
Block or sort encoding & $O_M(b_*^{12})$ vertices with leaves
 & $O_M(b_*^{16})$ vertices with color paths\\
Output matrix entries & $O_M(b_*^{24})$ & $O_M(b_*^{32})$\\
\bottomrule
\end{tabular}
\end{center}
The subcubic construction has $O_M(b_*^8)$ channels, with constantly
many auxiliary vertices per channel, and $O_M(b_*^8)$ sorts.
Replacing each auxiliary vertex by a gadget whose paths have length
at most the sort count gives the final $O_M(b_*^{16})$ order.
The proofs below justify these counts
and then account for the bit cost of producing the explicit matrices.
The constants may depend on the fixed $M,p$; the exponents in the
graph-size and compiler-time bounds are absolute.

\subsection{The smaller general compiler}

\begin{proposition}[Uniform general-graph compiler]\label{prop:uniform-general}
Fix the machine $M$ and polynomial $p$ used in
Proposition~\ref{prop:computation}. For a word $w$ of length $m$, put
$r=2(p(m)+2)$. The constructions of
\Cref{prop:computation,prop:consistency,prop:general-realization}
produce two equally large simple uncolored
graphs of order $O_M((r+1)^{12})$. Their explicit adjacency matrices and the
binary parameter $k=r$ can be produced by a deterministic multitape Turing
machine in time $O_{M,p}((r+2)^{80})$.
\end{proposition}
\begin{proof}
Since $p(m)\ge m+2$, we have $m\le r$. The number of original consistency
types and schemas is therefore $O_M((r+1)^2)$, and every vector dimension is
at most two. Each schema has $r$ constant-size binary coordinate entries.
Subdividing each wire with $4(r+1)$ internal types gives
$O_M((r+1)^3)$ types and links. Even copying an entire $r$-coordinate table
onto every link requires only $O_M((r+1)^4)$ entries.

For each type or link and each row there are at most four blocks: one chooses
two binary labels, or two label matches each determined by a binary label.
A site has at most four legal valuations. A square has two vector readouts
of dimensions at most two and two scalar readouts, hence at most $2^6$
valuations, a deliberately loose bound. Thus both the number $J$ of blocks
and the common number $N$ of base vertices are $O_M((r+1)^4)$; moreover
$J\le N$ because every block has a legal valuation in each version.
There are at most $J(N+1)$ leaves at any base vertex, giving total order
\begin{equation}\label{eq:graph-order-bound}
 N+NJ(N+1)\le N+N^2(N+1)=O_M((r+1)^{12}).
\end{equation}
The two output matrices have $O_M((r+1)^{24})$ entries in total.
There is at least one scalar real type, which alone supplies four sites per
row and two valuations per site. The graph order is at least $8r$, so the
output parameter is below the order.

Here is an explicit polynomial-time implementation. All lists are stored
sequentially, so no unit-cost random access is assumed.
\begin{enumerate}[label=\arabic*.,leftmargin=*]
\item Count $m$, evaluate the fixed polynomial $p$, and set $d,r,K$.
Enumerate the fixed machine alphabet, carry and borrow cases, initial
positions, and final-time tests. Fill each product table coordinate by
coordinate. Only the $\max\{1,m\}$ exceptional initial positions and the
$d+1$ possible blank suffix domains are generated. Apply the buffer and
dummy construction and enumerate wire and link identifiers.
\item Enumerate the constant number of local label choices for each
type-row or link-row pair. Record the named readouts of each block and
enumerate all of its local valuations. Test its affine or homogeneous
equation directly. Count $N$ and enumerate leaves by an owner identifier and
a local index. This produces vertex lists in both graph versions.
\item Enumerate ordered pairs of vertices and stream the corresponding
matrix bits. For base vertices, compare the constant number of named
readouts and their values, and check that the blocks differ. For a pair
involving a leaf, test owner incidence. Leaf--leaf pairs are nonadjacent.
The diagonal is zero. Finally write $r$ in binary.
\end{enumerate}

For a concrete bit-cost allowance, these descriptions use
$O_{M,p}((r+2)^{12})$ records, each of at most
$O_{M,p}((r+2)^2)$ bits; this includes full coordinate tables, identifiers,
valuation lists and counters. Allow the larger total work-space length
\[
 D=O_{M,p}((r+2)^{16}).
\]
The stated loops need at most $O_{M,p}((r+2)^{24})$ record-generation trials
and matrix entries. Each uses at most $O_{M,p}((r+2)^2)$ elementary lookups,
copies, comparisons or arithmetic operations; constructing a coordinate
table can scan all $r$ entries. A sequential list lookup or an elementary
binary arithmetic operation on at most $D$ bits costs $O(D^2)$ time by
ordinary scanning and schoolbook algorithms. The total allowance is thus
$O_{M,p}((r+2)^{24+2+32})$, which is within the asserted exponent~80.

Polynomial evaluation also fits this budget even if $p$ has signed
coefficients and cancellations. If $p(x)=\sum_{i=0}^q c_i x^i$ and
$A_p=1+\sum_i|c_i|$, every Horner intermediate at $m\ge0$ has absolute value
at most $A_p(m+1)^q$. Its bit length is $O_p(\log(r+2))$, and the number
$q$ of arithmetic stages is fixed. The case $m=0$ has $p(0)\ge2$ and causes
no exception. Reading the input and writing the $O(\log r)$-bit dimension
are included. At no stage is the list of $2^r$ full addresses, the
$2^d$ computation times, a consistency witness, or an offset family
enumerated.
\end{proof}


\subsection{The connected subcubic compiler}

\begin{proposition}[Uniform construction]
\label{prop:uniform-construction}
Fix a machine \(M\) and polynomial \(p\) satisfying
\Cref{prop:computation}, and apply the buffer construction
of \Cref{prop:consistency}.  On an input
word of length $m$, put
\[
 r=2(p(m)+2),\qquad K=r+1,\qquad b_*=r+2.
\]
The two graphs $X^0,X^1$ and the binary dimension $k=r$ can be produced
deterministically in time $O_{M,p}(b_*^{160})$.  Each output graph has
$O_M(b_*^{16})$ vertices.  They have the same positive order and are
connected, simple, uncolored, undirected graphs of maximum degree at most
three, regardless of whether the scalar-image instance is consistent.
\end{proposition}

\begin{proof}
We give explicit polynomial bounds and a construction that uses only the
compressed descriptions.  Since $p(m)\geq m+2$, we have $r\geq 8$ and
$m\leq b_*$.  The computation construction produces
$O_M((r+1)^2+m)=O_M(b_*^2)$ real types and comparison schemas.  Their
value dimensions are at most two.  A schema is described by its endpoint
identifiers, its two constant-size affine forms, and $r$ constant-size
coordinate domains and maps.

Each comparison wire has $4K$ private internal types and $4K+1$ links.
Thus the subdivided layout has $O_M(b_*^3)$ types and links.  At a fixed
row and type there are at most four choices of a site's pair of labels;
at a fixed row and link there are at most four choices of a square's two
coordinate matches.  Multiplying by $r$ rows gives
\begin{equation}
\label{uniform:block-bound}
 \abs{\cB}=O_M(b_*^4).
\end{equation}
A block has at most four readout names.  The total number of proper
channels, indexed by ordered pairs of distinct blocks and a shared name,
is therefore at most $4\abs{\cB}^2$.  Including channel zero at every
block, the total number of channels is $O_M(b_*^8)$.

Every legal valuation uses at most six bits by
\Cref{lem:block-fibers}.  Hence a channel has at most $2^6$ spine
vertices, at most seven prefix levels, and at most $2^6$ vertices on any
one of those levels.  These are absolute constants, independent of $r$.
There are also only constantly many vertex sorts per channel and one
anchor and anchor sort per block.  Consequently,
\begin{equation}
\label{uniform:auxiliary-bound}
 \abs{V(Y^\eps)}=O_M(b_*^8),\qquad
 s=O_M(b_*^8),
\end{equation}
where $s$ is the number of positive integer sort colors.

A colored vertex of color $c$ is replaced by three triangle vertices,
three connectors, and three private paths with $c$ new vertices each.
The replacement therefore has exactly $6+3c$ vertices.  Since
$1\leq c\leq s$, \eqref{uniform:auxiliary-bound} gives
\begin{equation}
\label{uniform:graph-bound}
 \abs{V(X^\eps)}
 \leq (6+3s)\abs{V(Y^\eps)}=O_M(b_*^{16}).
\end{equation}
The two explicit output matrices have $O_M(b_*^{32})$ entries, and the
output dimension $r$ has $O(\log b_*)$ binary bits.

All structural promises follow from \Cref{prop:subcubic-realization}.
In particular, equal order does not require a compatible global offset:
each individual block has a nonempty affine valuation space, and any of
its legal offsets supplies a bijection between its two block parts.
Their separate cardinalities agree and can be added over the partition.
There is at least one site block and a legal valuation in each version,
so the common graph order is positive.  Degree-one and degree-two vertices
are permitted because the required bound is at most three.

For completeness, the construction can be implemented with the following
finite lists and loops.
\begin{enumerate}
\item Evaluate the fixed polynomial $p$ and generate the circuit's schema
      domains and coordinate maps by the carry, borrow, and initialization
      cases of \Cref{prop:computation}.  Expand each connection into
      its comparison schemas, subdivide each comparison wire, and
      enumerate the type, link, and block lists.  A coordinate scan has
      length $r$; none of these steps ranges over the full set of
      addresses.
\item Compare the at most four names of each ordered pair of blocks to
      enumerate channels and their opposite-channel identifiers.  For
      every channel, enumerate bit strings of length at most six to list
      its legal valuations and occurring prefixes.  A prefix occurs if
      one of the constantly many full strings extending it satisfies the
      local equation.  Enumerate the sorts to assign their positive
      color numbers.  Fix all ordering choices by the same nested
      enumeration rules in the two versions.
\item Form the auxiliary vertex and edge lists from the defining rules.
      For example, enumerate unordered pairs of auxiliary vertices and
      test those rules.  Assign a vertex's incident ports in generated
      edge order; the list has length at most three.  Generate the
      triangle, connector, and private-path vertices using the sort
      numbers as path-length counters.
\item Enumerate pairs of output vertices and write their adjacency bits.
      A vertex record contains its original auxiliary vertex, its
      triangle/connector/path kind, its port, and, for a path vertex,
      its distance from the connector.  With the old edge and port lists,
      the replacement rules decide adjacency from such records.
\end{enumerate}
All local valuation tests have constant dimension.  The witness sets and
the families $\cH_U$ are used only in the correctness proof and are never
constructed by this algorithm.

Here is a deliberately loose bit-time accounting on a deterministic
multitape machine.  Individual records, coordinate tables, and counters
can be padded to $O_{M,p}(b_*^2)$ bits.  Identifiers and counters need only
$O_{M,p}(\log b_*)$ bits, and individual schema tables use $O(r)$ entries;
even a time or position address has only $d=r/2$ bits.  There are at most
$O_{M,p}(b_*^{32})$ records if the output matrix bits are included, so
$O_{M,p}(b_*^{40})$ bits suffice for lists and working storage, including
scratch space.  The fixed polynomial $p$ is evaluated on polynomially
bounded integers by ordinary integer arithmetic.

The ordered block-pair loop has $O_M(b_*^8)$ iterations, the auxiliary
vertex-pair loop has $O_M(b_*^{16})$, and the output pair loop has
$O_M(b_*^{32})$.  The other enumeration loops have no larger count.
The work of each iteration can be expressed using
$O_{M,p}(b_*^2)$ basic record operations: identifier lookup or update,
copying, comparison, and elementary integer arithmetic.  In particular,
edge information can be stored as endpoint and port records, without
copying whole adjacency rows into vertex records.  The total number of
these operations is bounded by $O_{M,p}(b_*^{40})$.

Implement an identifier lookup by a sequential scan of the stored lists;
all basic record operations can conservatively be bounded by the square
of the allowed tape length, using sequential scans and schoolbook
arithmetic.  Thus the preceding bounds give time at most
$O_{M,p}(b_*^{120})$, which is within the stated
$O_{M,p}(b_*^{160})$ allowance.  No constant or threshold in this
implementation depends on a separately fixed value of the growing
dimension.  Finally, $b_*$ is polynomial in $m$ for the fixed polynomial
$p$, so this is deterministic polynomial time in the source input length.
\end{proof}


\subsection{The two completeness reductions}

\begin{proof}[Proof of \Cref{thm:main}]
Membership for both languages is \Cref{prop:membership}, including
invalid encodings and dimensions larger than the graph order.

For hardness, fix any binary language \(L_0\in\EXPTIME\).
Choose a normalized deterministic decider \(M\) for its complement
and a fixed polynomial \(p\) as in \Cref{prop:computation-images}.
On \(w\) of length \(m\), generate the product circuit and its
scalar-image instance at \(r=2(p(m)+2)\ge8\). For \(\WL\), use
the general realization; for \(\SubWL\), use the connected subcubic
realization. Output the two explicit adjacency matrices and \(k=r\)
in binary. The compiler bounds are, respectively,
\(O_{M,p}((r+2)^{80})\) and \(O_{M,p}((r+2)^{160})\).
Since \(p\) is fixed, both are polynomial in \(m\); their polynomial
degrees may depend on the source language, as permitted for a
many-one reduction.

All required graph promises hold on every input, whether accepted
or rejected by \(M\). The two realization propositions verify the
hypotheses of \Cref{prop:game-reduction}. Together with the computation
reduction, this gives, for either output pair \(W^0,W^1\),
\[
\begin{aligned}
 w\in L_0
 &\Longleftrightarrow M\text{ does not accept }w\\
 &\Longleftrightarrow \text{the circuit has no true terminal test}\\
 &\Longleftrightarrow \text{the scalar-image instance is consistent}\\
 &\Longleftrightarrow W^0\equiv_r W^1.
\end{aligned}
\]
Thus each construction is a deterministic polynomial-time many-one
reduction with the required polarity. Because \(L_0\) was arbitrary,
both target languages are \(\EXPTIME\)-hard. No dimension-dependent
threshold or fixed-dimension lower bound supplies a step of these
uniform reductions.
\end{proof}


\begin{corollary}[Unary dimension]\label{cor:unary-dimension}
If the supplied dimension is encoded in unary instead of binary,
both \(\WL\) and \(\SubWL\) remain \(\EXPTIME\)-complete under
deterministic polynomial-time many-one reductions.
\end{corollary}

\begin{proof}
For membership, convert the unary dimension to binary and apply
\Cref{prop:membership}; the conversion is polynomial in the input length.
For hardness, the reductions just proved output
\(k=r=2(p(m)+2)\), which is polynomial in the source input length \(m\).
Writing this dimension in unary therefore preserves their polynomial
running time. The output graphs and the equivalence condition are
unchanged.
\end{proof}
